\section{Proposed Taxonomy}

To classify the corpus frameworks uniformly, we propose a taxonomy of six dimensions: specification, context, roles, execution, validation, and portability. These dimensions were chosen because they recur, under different names, across the studies analyzed, and they correspond to the structured fields extracted from each study. The taxonomy is presented here as a reference; its empirical validation against the corpus, including the strong dimensions and the weak ones, is reported in the Results section. Table~\ref{tab:taxonomia} summarizes each dimension through a guiding question and observable indicators.

\noindent\textbf{Specification.} Describes how the framework turns human intent into artifacts that guide the agent's work. At the weakest extreme, the specification is just a prompt; at the strongest, it is a versioned set of requirements, acceptance criteria, plans, tasks, architecture, and policies. Spec-driven development makes this dimension explicit by treating the specification as a versionable contract that drives generation \cite{DeepakBabuPiskala2026}, and AgileGen operationalizes it through executable acceptance criteria \cite{SaiZhang2025}. Reversa shifts the origin of the specification: instead of starting from a new product, it recovers operational specifications from existing systems \cite{Macedo2026}.

\noindent\textbf{Context.} The mechanism by which the framework decides what the agent should know before acting, which includes documents, code, history, architectural decisions, local rules, and collected evidence. Context engineering is treated as a discipline of its own for agents operating over real repositories \cite{SeyedmoeinMohsenimofidi2025}, and AutoCodeRover shows that structural retrieval of context from the code improves the quality of the patches generated \cite{YuntongZhang2024}. In Reversa, context is extracted from the legacy code and labeled by confidence \cite{Macedo2026}.

\noindent\textbf{Roles.} Define specialization, authority, and behavioral expectations. Multi-agent frameworks distribute tasks across personas such as product manager, architect, developer, tester, and reviewer, which reduces prompt ambiguity and sets output expectations. MetaGPT is the most explicit case, internalizing standard operating procedures and assigning specific artifacts to each role \cite{SiruiHong2023}, while ChatDev and CodePori show that role decomposition was already an active research line \cite{ChenQian2023,ZeeshanRasheed2024}.

\noindent\textbf{Execution.} Indicates whether the framework only produces guidance or also drives tools, edits code, runs tests, and interacts with interfaces. Modern agentic tools make execution a central dimension: the agent does not merely recommend, it modifies the environment. AutoDev and AutoCodeRover operate directly on the repository, editing code and running tests on real program-improvement tasks \cite{MicheleTufano2024,YuntongZhang2024}, and Prompt Sapper materializes execution as a production platform for AI chains \cite{YuCheng2023}.

\noindent\textbf{Validation.} Comprises the mechanisms for checking whether what was produced is correct, complete, and aligned with the context, which can include automated tests, human review, evidence artifacts, auditable traces, and readiness gates. Trace-based assurance articulates contracts, tests, and governance \cite{Paduraru2026}; verifiable collaboration among assistants makes every act in the flow auditable \cite{Merchant2025}; and Reversa makes confidence and gaps explicit in the recovered artifacts \cite{Macedo2026}. Validation is critical because agents can produce plausible results without real adherence to the system.

\noindent\textbf{Portability.} Describes whether the framework depends on a specific vendor, IDE, model, or format. In the corpus, this is the most fragile dimension: it appears almost always on the negative side, as a risk of lock-in and of dependence on a platform, and lock-in is named explicitly in only one study \cite{SeyedmoeinMohsenimofidi2025}. Almost no framework treats portability as a positive, deliberate design property, which is why it is best read as a predominantly diagnostic dimension, as detailed in the Results section.

\begin{table}[h!]
\centering
\caption{Dimensions of the proposed taxonomy.}
\label{tab:taxonomia}
\small
\begin{tabularx}{\textwidth}{p{2.8cm}Y Y}
\toprule
\textbf{Dimension} & \textbf{Guiding question} & \textbf{Observable indicators} \\
\midrule
Specification & How does intent become a work contract? & Specs, PRDs, plans, stories, tasks, acceptance criteria \\
Context & How does the agent know what is relevant? & Repository, docs, hooks, rules, memory, evidence, gaps \\
Roles & Who decides, who implements, who reviews? & Personas, agents, skills, responsibilities, authority \\
Execution & Does the framework act on the environment or only guide? & Code editing, commands, tests, browser, terminal, IDE \\
Validation & How are errors caught before they become deliverables? & Tests, checklists, gates, artifacts, human review, confidence \\
Portability & Does the process survive outside one tool? & Multiple integrations, open formats, lock-in, local install \\
\bottomrule
\end{tabularx}
\end{table}
